\documentclass[11pt, a4paper]{article}

\usepackage[utf8]{inputenc}
\usepackage[T1]{fontenc}
\usepackage[spanish]{babel}
\usepackage[top=2.5cm, bottom=2.5cm, left=2.5cm, right=2.5cm]{geometry}

% Fuente Palatino — más elegante y formal que Latin Modern
\usepackage{mathpazo}
\linespread{1.05}

\usepackage{xcolor}
\usepackage{enumitem}
\usepackage{titlesec}
\usepackage{microtype}
\usepackage{hyperref}
\usepackage{fancyhdr}
\usepackage{parskip}
\usepackage{mdframed}
\usepackage{booktabs}
\usepackage{tabularx}
\usepackage{colortbl}
\usepackage{array}

% ── Colores ──────────────────────────────────────────────────────────────────
\definecolor{navy}{HTML}{0f172a}
\definecolor{accent}{HTML}{1d4ed8}
\definecolor{okgreen}{HTML}{15803d}
\definecolor{lightrow}{HTML}{f8fafc}
\definecolor{bordergray}{HTML}{cbd5e1}
\definecolor{textsub}{HTML}{64748b}
\definecolor{headerbg}{HTML}{0f172a}
\definecolor{nextbg}{HTML}{eff6ff}
\definecolor{nextborder}{HTML}{93c5fd}

\hypersetup{colorlinks=true, urlcolor=accent, linkcolor=navy}

\pagestyle{fancy}
\fancyhf{}
\renewcommand{\headrulewidth}{0.4pt}
\fancyhead[L]{\small\color{textsub} DefyVision Metalconf}
\fancyhead[R]{\small\color{textsub} Reporte de avances — junio 2026}
\fancyfoot[C]{\small\color{textsub}\thepage}

\titleformat{\section}{\large\bfseries\color{navy}}{}{0em}{}[{\color{bordergray}\titlerule[0.6pt]}]
\titlespacing{\section}{0pt}{20pt}{8pt}

\setlist[itemize]{topsep=3pt, itemsep=4pt, leftmargin=1.6em}
\setlist[enumerate]{topsep=3pt, itemsep=4pt, leftmargin=1.8em}

\newmdenv[
  backgroundcolor=white,
  linecolor=bordergray,
  linewidth=0.7pt,
  roundcorner=4pt,
  innertopmargin=10pt,
  innerbottommargin=10pt,
  innerleftmargin=14pt,
  innerrightmargin=14pt,
  skipabove=5pt,
  skipbelow=5pt,
]{abox}

\newmdenv[
  backgroundcolor=nextbg,
  linecolor=nextborder,
  linewidth=1pt,
  roundcorner=4pt,
  innertopmargin=10pt,
  innerbottommargin=10pt,
  innerleftmargin=14pt,
  innerrightmargin=14pt,
  skipabove=5pt,
  skipbelow=5pt,
]{nbox}

\newcolumntype{C}[1]{>{\centering\arraybackslash}p{#1}}
\renewcommand{\arraystretch}{1.45}

% ════════════════════════════════════════════════════════════════════════════
\begin{document}
\thispagestyle{fancy}

\begin{center}
  {\fontsize{20}{24}\selectfont\bfseries\color{navy} DefyVision Metalconf}\\[6pt]
  {\large\color{textsub} Reporte de avances — semana del 9 al 12 de junio de 2026}\\[4pt]
\end{center}

{\color{navy}\rule{\textwidth}{1.5pt}}

\vspace{4pt}

% ── Estado ───────────────────────────────────────────────────────────────────
\section{Estado del sistema}

El sistema de inspección automática está funcionando correctamente con las
\textbf{dos cámaras activas} en planta. Durante esta semana se realizaron pruebas
exhaustivas y ajustes en ambas, logrando que el sistema reconozca correctamente
las piezas buenas del patrón \textbf{Microperforado} sin generar falsas alarmas.

\vspace{8pt}

\begin{tabularx}{\textwidth}{>{\bfseries}p{3.5cm} X C{5cm}}
  \arrayrulecolor{bordergray}
  \rowcolor{headerbg}
  \color{white}\bfseries Cámara &
  \color{white}\bfseries Descripción &
  \color{white}\bfseries Estado \\
  \hline
  \rowcolor{lightrow}
  Cámara 1 & Microperforado — lado izquierdo & Correcto, a testear por operarios. \\
  Cámara 2 & Microperforado — lado derecho   & Correcto, a testear por operarios. \\
  \hline
\end{tabularx}

% ── Avances ──────────────────────────────────────────────────────────────────
\section{Avances de esta semana}

\begin{abox}
\textbf{Calibración de cada cámara por separado}\\[4pt]
Cada cámara fue calibrada de forma independiente con imágenes reales tomadas en
planta. Así, los parámetros de una no interfieren con los de la otra, y cada una
reconoce correctamente la chapa que le corresponde inspeccionar.
\end{abox}

\begin{abox}
\textbf{Nuevo mapa de referencia de agujeros}\\[4pt]
Se generó un nuevo mapa con la disposición exacta de los agujeros que debe tener
una pieza correcta. Este mapa es la base con la que el sistema compara cada chapa
que pasa por la cámara. Ahora incluye solo los agujeros de la zona más estable de
la chapa, descartando los bordes donde la variación natural del material generaba
confusión.
\end{abox}

\begin{abox}
\textbf{Eliminación de falsas alarmas en Cámara 2}\\[4pt]
La Cámara 2 estaba detectando agujeros que no existían, causados por reflejos de
luz en el borde de la chapa. Se ajustó la forma en que el sistema interpreta la
imagen en esa zona, y el problema quedó resuelto.
\end{abox}

\begin{abox}
\textbf{Pantalla de alerta más clara para el operario}\\[4pt]
Cuando el sistema detecta una pieza defectuosa y detiene la máquina, la pantalla
ahora indica de forma grande y visible \textbf{en cuál de las dos cámaras se
detectó el problema}, para que el operario sepa exactamente a dónde dirigirse sin
necesidad de interpretación.
\end{abox}

\begin{abox}
\textbf{Historial de funcionamiento}\\[4pt]
Se incorporó una pantalla de historial donde se puede ver el desempeño del sistema
a lo largo del tiempo en cada cámara, con gráficos simples. Antes esta pantalla
aparecía vacía aunque hubiera datos guardados; eso quedó corregido.
\end{abox}

\begin{abox}
\textbf{Corrección automática de posición}\\[4pt]
Se desarrolló un mecanismo que detecta si la chapa se fue corriendo lentamente de
posición a lo largo del turno y ajusta el área de inspección de forma gradual y
silenciosa, sin interrumpir la producción. Este ajuste funciona de manera muy
conservadora: solo actúa cuando el desplazamiento es claro y sostenido durante
varios minutos.
\end{abox}

\begin{abox}
\textbf{Indicador de estado del área de inspección}\\[4pt]
En la pantalla de configuración ahora se muestra si el área que la cámara está
inspeccionando coincide correctamente con la zona de la chapa, o si se detecta
algún desplazamiento que requiera atención.
\end{abox}

\begin{abox}
\textbf{Mejor organización de grabaciones}\\[4pt]
Las carpetas donde se guardan las imágenes capturadas ahora incluyen en el nombre
el número de cámara correspondiente, haciendo más fácil identificar de qué estación
proviene cada lote.
\end{abox}

% ── Próximos pasos ────────────────────────────────────────────────────────────
\section{Próximos pasos}

\begin{nbox}
\begin{enumerate}
  \item Poner el sistema en marcha con ambas cámaras funcionando al mismo tiempo en
        condiciones reales de producción.
  \item Probar el sistema con piezas defectuosas reales para confirmar que la
        detección funciona correctamente.
  \item Seguir observando el comportamiento de la Cámara 2 con distintos lotes de
        material para confirmar la estabilidad de la calibración.
\end{enumerate}
\end{nbox}

% ── Pie ───────────────────────────────────────────────────────────────────────
\vspace{16pt}
{\color{bordergray}\rule{\textwidth}{0.6pt}}\\[4pt]
{\small\color{textsub}
  DefyVision\quad\textbullet\quad
  \href{mailto:design@defymotion.com.ar}{design@defymotion.com.ar}\quad\textbullet\quad
  Junio 2026
}

\end{document}
